\booksection{第二十章：開發者模式與高階許可權}

\chapterintro{這一章是給那些想要深入探索手機“隱藏功能”的使用者準備的。開發者模式、Root許可權、BL鎖……這些聽起來很極客的概念，其實離普通使用者並不遙遠。}

\booksubsection{一、開發者模式：解鎖隱藏設定}

開發者模式是手機系統內建的一個隱藏選單，裡面有很多普通使用者用不到但開發者需要的功能。

\subsubsection*{如何開啟：}

\begin{itemize}
  \item \textbf{安卓}：進入設定 > 關於手機 > 連續點選"版本號"7次，直到提示"您已處於開發者模式"。
  \item \textbf{鴻蒙}：設定 > 關於手機 > 連續點選"HarmonyOS版本"7次。
  \item \textbf{iOS}：普通使用者無法開啟完整的開發者模式，需要付費的Apple Developer賬號。
\end{itemize}

\subsubsection*{常用功能一覽：}

\begin{itemize}
  \item \textbf{USB除錯}：連線電腦進行高階操作（如刷機、資料恢復、ADB命令）的必開開關，詳見本章"ADB實用命令"小節。
  \item \textbf{顯示觸控操作}：螢幕上顯示觸控點（一個小圓點跟隨手指），適合錄製操作教程或演示。
\end{itemize}

\subsubsection*{效能與流暢度相關：}

\begin{itemize}
  \item \textbf{動畫速度調節（視窗動畫縮放 / 過渡動畫縮放 / 動畫程式時長縮放）}：這三個選項控制開啟App、切換頁面時的動畫時長。預設是1x，調到0.5x會讓手機"看起來"更快（實際上程式執行速度沒變，只是動畫變快了），調到關閉則直接無動畫跳轉。多數極客會設為0.5x，既保留過渡感又顯快；完全關閉會讓操作顯得生硬，並且可能掩蓋程式真實的卡頓，反而讓你誤以為"很流暢"。
  \item \textbf{強制GPU渲染（Force GPU Rendering）}：強制所有2D繪圖透過GPU而非CPU處理。在老舊機型或某些App最佳化不佳時能提升滾動流暢度，但也會增加GPU負載和耗電。現代手機預設已大量使用GPU渲染，這個開關對中高階機型影響有限，更多是歷史遺留選項，普通使用者保持預設即可。
  \item \textbf{顯示GPU呈現模式分析（Profile GPU Rendering）}：在螢幕上以柱狀圖實時顯示每幀的渲染耗時，綠色橫線代表16.6ms（60fps的幀時間閾值）。柱子超過綠線就說明掉幀了。這是開發者排查App卡頓的利器，普通使用者也可以用它來對比哪個App最佳化得差——比如某個App滾動時柱子經常飆到紅線以上，那基本就是它的鍋。
\end{itemize}

\subsubsection*{後臺與續航相關：}

\begin{itemize}
  \item \textbf{後台程式限制}：可以設定"標準限制""最多1個程式""最多2個程式"等。表面上看限制後台程式能省電，但實際上Android自帶的後臺管理已經做得不錯，強行限制反而會導致App頻繁重新載入（冷啟動反而更耗電、體驗也變差）。這個選項更適合除錯用，不建議普通使用者長期開啟，保持"標準限制"即可。
  \item \textbf{不保留活動}：每離開一個介面就立即銷燬它。極不推薦，會導致App頻繁白屏、重新載入，僅用於測試App的恢復邏輯。
\end{itemize}

\begin{tcbwarn}{注意事項}
開發者模式裡的設定可能會影響系統穩定性，不懂的選項不要亂改。開啟USB除錯後，連線電腦時要小心，不要連線不信任的電腦。

\end{tcbwarn}

\booksubsection{二、OEM解鎖：開發者選項裡的“刷機開關”}

在開發者選項裡有一個不太起眼但很重要的開關——“OEM解鎖”（部分品牌叫“OEM解鎖允許”或“裝置解鎖”）。這個開關本身不會解鎖BL鎖，但它是一個“前置許可”：只有開啟它，廠商才允許你透過fastboot命令或品牌工具去解鎖BL。

\begin{itemize}
  \item \textbf{開關的作用}：相當於向廠商宣告"我同意放棄保修、自負風險，請允許我解鎖BL"。如果這個開關是關閉的，任何解鎖BL的命令都會被直接拒絕。
  \item \textbf{與BL鎖的關係}：OEM解鎖是"申請資格"，BL解鎖才是"實際操作"。前者在手機設定裡點一下就行，後者需要去廠商官網申請解鎖碼、用電腦跑fastboot命令（如\texttt{fastboot flashing unlock}），二者缺一不可。
  \item \textbf{恢復出廠設定會重置嗎}：會。"OEM解鎖"開關在恢復出廠後會回到關閉狀態，但BL鎖一旦解鎖就是物理級別的狀態，不會因為恢復出廠而重新上鎖——需要手動執行\texttt{fastboot flashing lock}才能重新上鎖。
  \item \textbf{什麼時候應該開啟}：除非你確定要刷機或Root，否則保持關閉。它和"USB除錯"是兩個獨立開關，刷機時通常需要同時開啟。
\end{itemize}

\begin{tcbinfo}{小知識：為什麼廠商要設定這個開關}
早期安卓沒有這個開關，惡意軟體只要拿到許可權就能偷偷刷入系統。Google後來要求所有安卓廠商在開發者選項中加入“OEM解鎖”開關，作為一道“使用者知情同意”的確認。這也就是為什麼即便你解了BL鎖，如果OEM解鎖開關被關閉，重新刷機仍然會失敗。

\end{tcbinfo}

\booksubsection{三、BL鎖（Bootloader鎖）：刷機的第一道門}

Bootloader是手機啟動時載入作業系統的程式。BL鎖就是廠商為了安全而加的鎖，防止使用者刷入非官方系統。

\subsubsection*{BL鎖的作用：}

\begin{itemize}
  \item \textbf{安全保護}：防止惡意軟體刷入系統。
  \item \textbf{防盜}：即使手機被刷機，沒有原賬號密碼也無法使用。
  \item \textbf{保修限制}：解鎖BL通常意味著失去官方保修。
\end{itemize}

\subsubsection*{各品牌BL鎖政策：}

\begin{itemize}
  \item \textbf{小米}：可以解鎖，需要繫結賬號等待期（7-30天）。
  \item \textbf{華為}：2018年後停止提供解鎖碼，基本無法官方解鎖。
  \item \textbf{一加}：可以解鎖，相對寬鬆。
  \item \textbf{OPPO/vivo}：部分機型可以解鎖，需要申請。
  \item \textbf{三星}：可以解鎖，但會觸發Knox熔斷，失去Samsung Pay等功能。
  \item \textbf{蘋果}：iOS裝置沒有傳統意義上的BL鎖概念，但有類似的安全機制。
\end{itemize}

\booksubsection{四、ADB實用命令：透過電腦控制手機}

ADB 全稱 Android Debug Bridge（安卓除錯橋），是 Google 官方提供的命令列工具，讓電腦透過 USB 或 Wi-Fi 與安卓手機通訊。它本身不需要 Root，只要手機開了“USB除錯”就能用。可以把它理解為“用電腦給手機下指令的遙控器”——刷機、裝App、傳檔案、查日誌，都得靠它。

\subsubsection*{如何透過 USB 連線：}

\begin{enumerate}
  \item 手機端：進入開發者選項 開啟"USB除錯"。
  \item 電腦端：從 Google 官方下載 Platform Tools（含\texttt{adb}和\texttt{fastboot}兩個可執行檔案），解壓到任意目錄。
  \item 用資料線連線手機和電腦，手機上會彈窗"是否允許USB除錯"，勾選"始終允許"並確定。
  \item 在電腦命令列（Windows 的 PowerShell 或 cmd）裡執行\texttt{adb devices}，如果輸出類似\texttt{List of devices attached /\allowbreak  xxxxxxxx device}，就說明連上了。
\end{enumerate}

\begin{tcbwarn}{驅動和授權彈窗別忽略}
Windows 使用者如果\texttt{adb devices}顯示 “unauthorized” 或根本不列出裝置，多半是驅動沒裝好或者忘了在手機上點“允許”。三星/小米等品牌可以裝各自的“USB驅動”或“手機助手”來補齊驅動；Mac 和 Linux 一般免驅。如果手機彈窗被你手滑點掉了，拔插一下資料線就會重新彈出來。

\end{tcbwarn}

\subsubsection*{常用 ADB 命令：}

\begin{itemize}
  \item \texttt{adb install app.\allowbreak apk}：把電腦上的 APK 裝到手機（\texttt{-\allowbreak r}表示覆蓋安裝並保留資料；\texttt{-\allowbreak s}裝到 SD 卡）。
  \item \texttt{adb uninstall com.\allowbreak xxx.\allowbreak xxx}：按包名解除安裝 App。
  \item \texttt{adb push 電腦路徑 /\allowbreak sdcard/\allowbreak 目標路徑}：把電腦檔案推送到手機。
  \item \texttt{adb pull /\allowbreak sdcard/\allowbreak 檔案 電腦路徑}：把手機檔案拉到電腦（比發給自己微信再下載快得多）。
  \item \texttt{adb shell}：進入手機的 Linux Shell，可以直接執行\texttt{ls}、\texttt{pm}、\texttt{am}等命令。
  \item \texttt{adb reboot bootloader}：重啟手機進入 Fastboot 模式（刷機的入口）。
  \item \texttt{adb reboot recovery}：重啟進入 Recovery 模式。
  \item \texttt{adb shell pm list packages}：列出手機上所有已安裝應用的包名（找包名必用）。加\texttt{-\allowbreak s}只看系統應用，加\texttt{-\allowbreak 3}只看第三方應用。
  \item \texttt{adb shell pm uninstall -\allowbreak -\allowbreak user 0 套件名稱}：為當前使用者解除安裝預裝 App，免 Root，詳見下方實戰。
\end{itemize}

\subsubsection*{免 Root 解除安裝預裝 App 實戰：}

這是 ADB 最實用的功能之一：用\texttt{pm uninstall -\allowbreak -\allowbreak user 0}命令可以為“當前使用者”解除安裝預裝 App，而不需要 Root，也不影響系統分割區。被解除安裝的 App 在系統裡“消失”，但實際還在系統分割區裡，恢復出廠設定後會回來。

\begin{enumerate}
  \item 先用\texttt{adb shell pm list packages -\allowbreak s}列出所有系統應用，找到你想解除安裝的包名。比如小米瀏覽器是\texttt{com.\allowbreak android.\allowbreak browser}，華為應用市場是\texttt{com.\allowbreak vmall.\allowbreak client}。
  \item 不確定包名對應哪個 App？可以在電腦上用\texttt{adb shell pm list packages | findstr 關鍵字}搜尋，或者裝一個第三方 App "Package Viewer"在手機上直接檢視每個 App 的包名。
  \item 執行解除安裝：\texttt{adb shell pm uninstall -\allowbreak -\allowbreak user 0 套件名稱}，比如\texttt{adb shell pm uninstall -\allowbreak -\allowbreak user 0 com.\allowbreak android.\allowbreak browser}。返回\texttt{Success}就成功了。
  \item 想恢復怎麼辦？兩種方法：(1) 恢復出廠設定（最徹底但麻煩）；(2) 重新安裝：\texttt{adb shell cmd package install-\allowbreak existing 套件名稱}。
\end{enumerate}

\textbf{千萬別亂刪系統核心包}\par

“pm uninstall --user 0”雖然是免 Root 的安全方式，但亂刪關鍵系統應用（如\texttt{com.\allowbreak android.\allowbreak systemui}狀態列、\texttt{com.\allowbreak android.\allowbreak phone}電話、\texttt{com.\allowbreak android.\allowbreak settings}設定）會導致系統崩潰、黑屏、無法開機等嚴重問題。刪之前先 Google 一下“這個包名能不能刪”，寧可少刪不要錯刪。建議一次只刪一兩個，確認沒問題再繼續。

\begin{tcbinfo}{小知識：--user 0 是什麼意思}
安卓 4.2 之後支援多使用者，每個使用者有獨立的 ID（機主是 0，附加使用者從 10 開始）。\texttt{-\allowbreak -\allowbreak user 0}表示“只對當前機主解除安裝”，系統分割區裡的 APK 檔案依然存在，只是不給機主顯示了。這就是為什麼它不需要 Root——沒動系統檔案，只是改了“使用者配置”。這也是它和 Root 後解除安裝最大的區別：Root 解除安裝是真正從系統裡刪掉，所有使用者都不再有；ADB 解除安裝只是“對當前使用者隱藏”。

\end{tcbinfo}

\booksubsection{五、Root許可權：安卓的“超級管理員”}

Root許可權相當於Linux中的超級使用者（對/目錄擁有完全控制權，可以讀寫任意檔案），是安卓系統的最高許可權。它與Windows的Administrator有本質區別：Windows出於安全考慮對管理員做了諸多限制（真正的最高許可權其實是 TrustedInstaller），即便是 Administrator 也無法像 Linux root 那樣直接執行\texttt{sudo rm -\allowbreak rf /\allowbreak *}。獲得Root權限後，你可以修改系統檔案、解除安裝預裝應用、使用需要Root的高階App。

\subsubsection*{Root的好處：}

\begin{itemize}
  \item 解除安裝無法解除安裝的預裝應用。
  \item 使用高階功能App（如鈦備份、綠色守護）。
  \item 修改系統介面和功能。
  \item 安裝第三方Recovery和自定義ROM。
\end{itemize}

\subsubsection*{Root的風險：}

\begin{itemize}
  \item \textbf{失去保修}：大多數廠商不保修Root過的手機。
  \item \textbf{安全風險}：惡意App可能獲得Root權限，造成嚴重後果。
  \item \textbf{系統不穩定}：誤操作可能導致系統崩潰。
  \item \textbf{銀行App不能用}：很多銀行和支付App檢測到Root會拒絕執行。
  \item \textbf{OTA更新失敗}：可能無法正常接收官方系統更新。
\end{itemize}

\textbf{嚴重警告}\par

Root操作有一定風險，可能導致手機變磚（無法開機）。如果你是新手，建議先充分了解相關知識，並確保有完整的備份。操作前請仔細閱讀教程，確認你的機型和系統版本。

\booksubsection{六、SafetyNet與Play Integrity認證：為什麼Root後銀行App不能用}

很多Root使用者都遇到過這樣的崩潰瞬間：Root之後興致勃勃地裝了一堆模組，結果開啟銀行App——“您的裝置不安全，無法使用”，Google Pay也用不了，連Netflix都拒絕高畫質播放。這背後是Google的兩套裝置完整性認證機制：舊的 SafetyNet 和新的 Play Integrity API。

\subsubsection*{什麼是 SafetyNet / Play Integrity：}

\begin{itemize}
  \item \textbf{SafetyNet Attestation（舊版，2017年起）}：Google 早期推出的裝置完整性檢查 API，應用透過它向 Google 伺服器查詢"這臺裝置是否透過 CTS 相容性測試、是否處於未修改的官方狀態"。返回結果分為 Basic Attestation（基礎校驗）和 Strong Attestation（強校驗，依賴硬體金鑰）。Google 已宣佈 SafetyNet 將在 2024 年後逐步停用。
  \item \textbf{Play Integrity API（新版，2022年起）}：Google 推出的替代方案，比 SafetyNet 更嚴格也更全面。除了裝置完整性，還包括"應用完整性"（App 是否被篡改）和"帳號/授權完整性"（使用者帳號是否被 Google 信任）。新 App 應遷移到 Play Integrity。
  \item \textbf{誰來用}：銀行、金融、支付類 App（如 Google Pay、各銀行的網銀、支付寶/微信支付的部分場景）、流媒體（Netflix 高畫質版權保護）、部分遊戲反作弊（如騰訊/網易的部分遊戲）、企業辦公 App（MDM 裝置合規性檢查）。
\end{itemize}

\subsubsection*{Play Integrity 的三級認證結果：}

\begin{itemize}
  \item \textbf{MEETS\_DEVICE\_INTEGRITY（裝置完整性透過）}：裝置執行正版 Android 系統，Bootloader 未解鎖，沒有 Root，CTS 相容性測試透過。可以正常使用銀行、支付類 App。
  \item \textbf{MEETS\_BASIC\_INTEGRITY（基礎完整性透過）}：裝置基本可信，但可能存在輕微異常。某些 App 會降級使用（如 Netflix 仍可播放但限制清晰度），某些嚴格的銀行 App 會直接拒絕。
  \item \textbf{MEETS\_STRONG\_INTEGRITY（強完整性透過）}：最高階別，要求硬體金鑰背書（如 Google Pixel 系列的 Titan M 晶片），大多數第三方品牌機型即使未 Root 也達不到這一級。
  \item \textbf{無任何標籤（NO\_INTEGRITY）}：裝置已 Root、解鎖 BL、刷了第三方 ROM 或被識別為模擬器，絕大多數支付/銀行 App 會直接拒絕執行。
\end{itemize}

\subsubsection*{為什麼 Root 後這些 App 會失效：}

應用透過 Play Integrity API 向 Google 伺服器發起認證請求，Google 會檢查裝置的多項指標——Bootloader 鎖定狀態、系統分割區是否被修改、是否存在 su 二進位制檔案、SELinux 是否處於寬鬆模式等。任何一項異常都會導致返回的認證級別降低或為空，銀行 App 一看“裝置不安全”，自然就拒絕服務了。

\subsubsection*{現代隱藏方案：Magisk Hide / Shamiko / Zygisk：}

為了在享受 Root 便利的同時還能用銀行 App，社群開發了一系列“隱藏 Root 狀態”的方案。核心思路都是：讓檢測程式看到的不是一個被修改過的系統，而是一個“看起來正常”的裝置。

\begin{itemize}
  \item \textbf{Magisk}：當前最主流的 Systemless Root 方案——不修改系統分割區，而是透過掛載的方式實現 Root。這樣設計的好處是 OTA 更新失敗的機率小、解除安裝比較乾淨，是後續所有隱藏方案的基礎。但使用 Magisk 仍然需要解鎖 BL 鎖。
  \item \textbf{Magisk Hide（已停更）}：Magisk 早期內建的隱藏模組，能對指定 App 隱藏 Root。但 2021 年起已被 Magisk 官方移除，原因是 Google 持續升級檢測手段，雙方陷入"貓鼠遊戲"。
  \item \textbf{Zygisk}：Magisk 24.0+ 引入的新機制，類似 Riru，透過注入 Zygote 程式實現模組載入，比直接修改系統更隱蔽。開啟 Zygisk 是使用現代隱藏方案的前提。
  \item \textbf{Shamiko}：當前最主流的隱藏模組，配合 Zygisk 工作，可以為指定 App（"反向白名單"模式）隱藏 Root、隱藏 Magisk 自身、偽造 Bootloader 鎖定狀態。Shamiko 透過 LSPosed 倉庫或 GitHub 分發。
  \item \textbf{Play Integrity Fix（osm0sis）}：專門針對 Google Play Integrity 檢測的 Magisk 模組，會修改裝置指紋（fingerprint）讓 Google 認為你的裝置是某個未 Root 的官方機型。
\end{itemize}

\begin{tcbwarn}{貓鼠遊戲永遠不會結束}
Google 每隔一段時間就會更新檢測規則，App 廠商也會跟進升級 Play Integrity 校驗。今天還能用的隱藏方案，明天可能就被識別。如果你完全依賴銀行/支付 App（比如日常轉帳、掃碼付款），最穩妥的選擇仍然是：不要 Root、不要解 BL 鎖。隱藏方案更適合“主力機不 Root + 備用機玩 Root”的玩家配置。

\end{tcbwarn}

\booksubsection{七、給普通使用者的建議}

\begin{itemize}
  \item \textbf{開發者模式}：可以開啟，體驗一下動畫加速、GPU 呈現分析等功能，但不要亂改其他不理解的選項。
  \item \textbf{OEM 解鎖與 BL 鎖}：如果你不打算刷機或 Root，兩個開關都保持預設（關閉）狀態最安全。
  \item \textbf{ADB}：免 Root 的強大工具，普通人也能用——比如用\texttt{pm uninstall -\allowbreak -\allowbreak user 0}解除安裝預裝 App。但刪之前一定要查清楚包名能不能刪。
  \item \textbf{Root}：普通使用者不建議 Root，風險大於收益。如果確實需要某些功能，可以先了解有沒有不需要 Root 的替代方案（比如用 ADB 解除安裝預裝、用 Shizuku 替代部分 Root 功能）。
\end{itemize}

\begin{tcbpractice}{來動手：開啟開發者模式並調整動畫速度}
\begin{enumerate}
  \item 進入設定 > 關於手機。
  \item 連續點選"版本號"7次，直到提示已開啟開發者模式。
  \item 返回設定，找到"開發者選項"（通常在設定底部或"系統"裡）。
  \item 找到"視窗動畫縮放"、"過渡動畫縮放"、"動畫程式時長縮放"，將它們都設為0.5x或關閉。
  \item 感受一下手機是否變快了。
\end{enumerate}

\end{tcbpractice}
